% Options for packages loaded elsewhere
\PassOptionsToPackage{unicode}{hyperref}
\PassOptionsToPackage{hyphens}{url}
%
\documentclass[
]{article}
\usepackage{amsmath,amssymb}
\usepackage{iftex}
\ifPDFTeX
  \usepackage[T1]{fontenc}
  \usepackage[utf8]{inputenc}
  \usepackage{textcomp} % provide euro and other symbols
\else % if luatex or xetex
  \usepackage{unicode-math} % this also loads fontspec
  \defaultfontfeatures{Scale=MatchLowercase}
  \defaultfontfeatures[\rmfamily]{Ligatures=TeX,Scale=1}
\fi
\usepackage{lmodern}
\ifPDFTeX\else
  % xetex/luatex font selection
\fi
% Use upquote if available, for straight quotes in verbatim environments
\IfFileExists{upquote.sty}{\usepackage{upquote}}{}
\IfFileExists{microtype.sty}{% use microtype if available
  \usepackage[]{microtype}
  \UseMicrotypeSet[protrusion]{basicmath} % disable protrusion for tt fonts
}{}
\makeatletter
\@ifundefined{KOMAClassName}{% if non-KOMA class
  \IfFileExists{parskip.sty}{%
    \usepackage{parskip}
  }{% else
    \setlength{\parindent}{0pt}
    \setlength{\parskip}{6pt plus 2pt minus 1pt}}
}{% if KOMA class
  \KOMAoptions{parskip=half}}
\makeatother
\usepackage{xcolor}
\usepackage[margin=1in]{geometry}
\usepackage{color}
\usepackage{fancyvrb}
\newcommand{\VerbBar}{|}
\newcommand{\VERB}{\Verb[commandchars=\\\{\}]}
\DefineVerbatimEnvironment{Highlighting}{Verbatim}{commandchars=\\\{\}}
% Add ',fontsize=\small' for more characters per line
\usepackage{framed}
\definecolor{shadecolor}{RGB}{248,248,248}
\newenvironment{Shaded}{\begin{snugshade}}{\end{snugshade}}
\newcommand{\AlertTok}[1]{\textcolor[rgb]{0.94,0.16,0.16}{#1}}
\newcommand{\AnnotationTok}[1]{\textcolor[rgb]{0.56,0.35,0.01}{\textbf{\textit{#1}}}}
\newcommand{\AttributeTok}[1]{\textcolor[rgb]{0.13,0.29,0.53}{#1}}
\newcommand{\BaseNTok}[1]{\textcolor[rgb]{0.00,0.00,0.81}{#1}}
\newcommand{\BuiltInTok}[1]{#1}
\newcommand{\CharTok}[1]{\textcolor[rgb]{0.31,0.60,0.02}{#1}}
\newcommand{\CommentTok}[1]{\textcolor[rgb]{0.56,0.35,0.01}{\textit{#1}}}
\newcommand{\CommentVarTok}[1]{\textcolor[rgb]{0.56,0.35,0.01}{\textbf{\textit{#1}}}}
\newcommand{\ConstantTok}[1]{\textcolor[rgb]{0.56,0.35,0.01}{#1}}
\newcommand{\ControlFlowTok}[1]{\textcolor[rgb]{0.13,0.29,0.53}{\textbf{#1}}}
\newcommand{\DataTypeTok}[1]{\textcolor[rgb]{0.13,0.29,0.53}{#1}}
\newcommand{\DecValTok}[1]{\textcolor[rgb]{0.00,0.00,0.81}{#1}}
\newcommand{\DocumentationTok}[1]{\textcolor[rgb]{0.56,0.35,0.01}{\textbf{\textit{#1}}}}
\newcommand{\ErrorTok}[1]{\textcolor[rgb]{0.64,0.00,0.00}{\textbf{#1}}}
\newcommand{\ExtensionTok}[1]{#1}
\newcommand{\FloatTok}[1]{\textcolor[rgb]{0.00,0.00,0.81}{#1}}
\newcommand{\FunctionTok}[1]{\textcolor[rgb]{0.13,0.29,0.53}{\textbf{#1}}}
\newcommand{\ImportTok}[1]{#1}
\newcommand{\InformationTok}[1]{\textcolor[rgb]{0.56,0.35,0.01}{\textbf{\textit{#1}}}}
\newcommand{\KeywordTok}[1]{\textcolor[rgb]{0.13,0.29,0.53}{\textbf{#1}}}
\newcommand{\NormalTok}[1]{#1}
\newcommand{\OperatorTok}[1]{\textcolor[rgb]{0.81,0.36,0.00}{\textbf{#1}}}
\newcommand{\OtherTok}[1]{\textcolor[rgb]{0.56,0.35,0.01}{#1}}
\newcommand{\PreprocessorTok}[1]{\textcolor[rgb]{0.56,0.35,0.01}{\textit{#1}}}
\newcommand{\RegionMarkerTok}[1]{#1}
\newcommand{\SpecialCharTok}[1]{\textcolor[rgb]{0.81,0.36,0.00}{\textbf{#1}}}
\newcommand{\SpecialStringTok}[1]{\textcolor[rgb]{0.31,0.60,0.02}{#1}}
\newcommand{\StringTok}[1]{\textcolor[rgb]{0.31,0.60,0.02}{#1}}
\newcommand{\VariableTok}[1]{\textcolor[rgb]{0.00,0.00,0.00}{#1}}
\newcommand{\VerbatimStringTok}[1]{\textcolor[rgb]{0.31,0.60,0.02}{#1}}
\newcommand{\WarningTok}[1]{\textcolor[rgb]{0.56,0.35,0.01}{\textbf{\textit{#1}}}}
\usepackage{graphicx}
\makeatletter
\newsavebox\pandoc@box
\newcommand*\pandocbounded[1]{% scales image to fit in text height/width
  \sbox\pandoc@box{#1}%
  \Gscale@div\@tempa{\textheight}{\dimexpr\ht\pandoc@box+\dp\pandoc@box\relax}%
  \Gscale@div\@tempb{\linewidth}{\wd\pandoc@box}%
  \ifdim\@tempb\p@<\@tempa\p@\let\@tempa\@tempb\fi% select the smaller of both
  \ifdim\@tempa\p@<\p@\scalebox{\@tempa}{\usebox\pandoc@box}%
  \else\usebox{\pandoc@box}%
  \fi%
}
% Set default figure placement to htbp
\def\fps@figure{htbp}
\makeatother
\setlength{\emergencystretch}{3em} % prevent overfull lines
\providecommand{\tightlist}{%
  \setlength{\itemsep}{0pt}\setlength{\parskip}{0pt}}
\setcounter{secnumdepth}{-\maxdimen} % remove section numbering
\usepackage{tikz}
\usetikzlibrary{arrows.meta, positioning}
\usepackage{bookmark}
\IfFileExists{xurl.sty}{\usepackage{xurl}}{} % add URL line breaks if available
\urlstyle{same}
\hypersetup{
  pdftitle={Informe TP 2},
  hidelinks,
  pdfcreator={LaTeX via pandoc}}

\title{Informe TP 2}
\usepackage{etoolbox}
\makeatletter
\providecommand{\subtitle}[1]{% add subtitle to \maketitle
  \apptocmd{\@title}{\par {\large #1 \par}}{}{}
}
\makeatother
\subtitle{Modelado Estadistico}
\author{\textbf{Integrantes del equipo}\\
- Alamo, Malena (/21)\\
- Corzini, Fabrizio (32/22)\\
- Laria, Jeremias (/21)}
\date{10 /07 /2025}

\begin{document}
\maketitle

\subsection{Consigna 1}\label{consigna-1}

\begin{quote}
\emph{a completar al final porque es opcional}
\end{quote}

\begin{center}\rule{0.5\linewidth}{0.5pt}\end{center}

\subsection{Consigna 2}\label{consigna-2}

\subsubsection{2.1}\label{section}

\begin{quote}
\emph{Plantear un modelo de efectos fijos para predecir el puntaje de
IMDB únicamente en función del paı́s de origen.}
\end{quote}

\subsubsection{2.2}\label{section-1}

\begin{quote}
\emph{Plantear un modelo de efectos aleatorios para predecir el puntaje
de IMDB únicamente en función del paı́s de origen.}
\end{quote}

\subsubsection{2.3}\label{section-2}

\begin{quote}
\emph{Mostrar las estimaciones de los efectos de ambos modelos en un
mismo gráfico e interpretar cómo se diferencian.}
\end{quote}

\begin{center}\rule{0.5\linewidth}{0.5pt}\end{center}

\subsection{Consigna 3}\label{consigna-3}

\begin{quote}
\emph{Usando únicamente la variable release year, predecir la
popularidad de cada tı́tulo (usando un tipo de modelo que crea adecuado)
con un spline cúbico. Usar k = 1, 2, 3, 5, 10, 20, 50 nodos fijando el λ
(penalización de rugosidad) en 0, y comparar todas las curvas estimadas
en un mismo gráfico.}
\end{quote}

\begin{center}\rule{0.5\linewidth}{0.5pt}\end{center}

\subsection{Consigna 4}\label{consigna-4}

\begin{quote}
\emph{A qué subconjunto de las variables Año, Duración y Paı́s de debe
condicionar para estimar el efecto causal promedio de la variable
Comedia sobre el Score? Dar todas las posibilidades.}
\end{quote}

Para poder estimar el efecto causal de una variable X por sobre Y hay
que primero eliminar todo tipo de asociación que pueda existir entre
estas variables (todo tipo de relación subyacente que pueda estar
ocurriendo por detrás que explique la variable Y que no se trate del
efecto de X).

Para eliminar esas asociaciones lo que se puede hacer es
\emph{condicionar} a X bajo un cierto subconjunto de variables que fije
el efecto de variables exógenas. Esto se hace de la siguiente manera:

\begin{enumerate}
\def\labelenumi{\arabic{enumi}.}
\tightlist
\item
  Copiamos el DAG original pero le sacamos las aristas que \emph{salen}
  de \textbf{Comedia}.
\item
  Identificamos todos los \emph{caminos} que existan entre
  \textbf{Comedia} y \textbf{Score}.
\item
  Construimos todos los posibles subconjuntos que garanticen una
  d-separación entre \textbf{Comedia} y \textbf{Score} (o análogamente
  que bloqueen todo camino entre ellos).
\end{enumerate}

\subsubsection{DAG}\label{dag}

\begin{Shaded}
\begin{Highlighting}[]
\KeywordTok{\textbackslash{}begin}\NormalTok{\{}\ExtensionTok{center}\NormalTok{\}}
\KeywordTok{\textbackslash{}begin}\NormalTok{\{}\ExtensionTok{tikzpicture}\NormalTok{\}[}
\NormalTok{  node distance=2cm,}
\NormalTok{  every node/.style=\{draw, ellipse, minimum width=2.5cm, align=center\},}
\NormalTok{  \textgreater{}=\{Stealth[length=3mm]\}}
\NormalTok{]}
\FunctionTok{\textbackslash{}node}\NormalTok{ (U) \{U\};}
\FunctionTok{\textbackslash{}node}\NormalTok{[right=of U] (X) \{X\};}
\FunctionTok{\textbackslash{}node}\NormalTok{[right=of X] (Y) \{Y\};}

\FunctionTok{\textbackslash{}draw}\NormalTok{[{-}\textgreater{}] (U) {-}{-} (X);}
\FunctionTok{\textbackslash{}draw}\NormalTok{[{-}\textgreater{}] (U) {-}{-} (Y);}
\FunctionTok{\textbackslash{}draw}\NormalTok{[{-}\textgreater{}] (X) {-}{-} (Y);}
\KeywordTok{\textbackslash{}end}\NormalTok{\{}\ExtensionTok{tikzpicture}\NormalTok{\}}
\KeywordTok{\textbackslash{}end}\NormalTok{\{}\ExtensionTok{center}\NormalTok{\}}
\end{Highlighting}
\end{Shaded}

\begin{center}\rule{0.5\linewidth}{0.5pt}\end{center}

\subsection{Consigna 5}\label{consigna-5}

\begin{quote}
\emph{Dividir al conjunto de datos en entrenamiento y testeo (tambien
puede usar otra técnica, como validación cruzada). Con todas las
variables que tiene disponibles, probar al menos 3 modelos diferentes y
elegir el que minimice el error cuadrático médio de predicción para el
rating de IMDB.}
\end{quote}

\begin{center}\rule{0.5\linewidth}{0.5pt}\end{center}

\subsection{Consigna 6}\label{consigna-6}

\begin{quote}
\emph{A partir del modelo elegido en el item anterior, producir un
archivo predicciones.csv que tenga una sola columna que contenga, en la
fila i, la predicción del rating de IMDB para el tı́tulo de la fila i.}
\end{quote}

\begin{center}\rule{0.5\linewidth}{0.5pt}\end{center}

\subsection{Referencias}\label{referencias}

\begin{itemize}
\tightlist
\item
  Clases Teoricas\\
\item
  Introduction to Causal Inference, Brady Neal
\end{itemize}

\end{document}
